% !TeX root = ../教材题库_全部条目_转写工作稿.tex
% Source: XJ-B1-C06-TH-P249-Y0892
% 第6章 / 6.4 用样本估计总体 / 教材 p.249 / PDF 第256页
% OCR未校对：公式与文字均可能有误；未完成数学TeX转写；禁止直接发给学生
% 原题及相邻原解截图（核对用）：../题目原图/XJ-B1-C06-TH-P249-Y0892-01.webp
% 本文件只在转写完成后作为题干使用；OCR 条目尚不是可靠题干。
\begin{bankitem}{p.249 思考（含紧邻说明）}{第6章 / 6.4 用样本估计总体 / 教材 p.249 / PDF 第256页}{XJ-B1-C06-TH-P249-Y0892}{OCR未校对：公式与文字均可能有误；未完成数学TeX转写；禁止直接发给学生}
\begin{ocrdraft}
[教材p.249 / PDF第256页]
\par
？思考例9中得到的平均数和中位数大小相等，这一关系具有普遍性吗？
\par
一般地，对一个单峰的频率分布直方图来说，如果直方图的形状是对称的（如图
\par
6．4一4（1)），那么平均数和中位数应该大体上差不多；如果直方图在右边矮长些（如
\par
图6．4一4（2）），那么平均数大于中位数；如果直方图在左边矮长些（如图6．4一4（3）），
\par
那么平均数小于中位数．也就是说，与中位数相比，平均数总在矮长的那一边．
\par
平均数、中位数
\par
中位数平均数
\par
平均数中位数
\par
（2）
\par
（3）
\par
图6．4一4
\par
结合6，3节知识可知，随着样本容量的不断增大，当频率分布直方图的组数不
\par
断增多，分组的组距不断缩小时，较多的原始数据信息被保留，从而频率分布直方
\par
图越来越接近于总体的分布．于是，我们可用频率分布直方图估计出的样本的平均
\par
数、中位数、众数等来作为总体的相应估计值．但要注意的是，用频率分布直方图
\par
比通过样本数据得来的估计值的精度低，且估计结果与数据分组有关．
\end{ocrdraft}
\end{bankitem}
